\section{Introducción}

Medir la participación de un egresado es, paradójicamente, más difícil cuanto más valioso resulta esa participación. La universidad acompaña al estudiante con un detalle minucioso a lo largo de su formación —cada materia, cada nota, cada semestre queda registrado—, pero pierde de vista al graduado justo en el momento en que comienza lo que mejor refleja el valor de aquello que aprendió: su inserción en el mercado laboral, la creación de una empresa, la producción de conocimiento científico. La trayectoria que más importa para juzgar la pertinencia de un programa es, precisamente, la que ocurre fuera de los muros de la institución y, por tanto, fuera de sus sistemas de información.

El obstáculo no es solo de seguimiento, sino de observación. Ninguna fuente única alcanza a mirar al egresado en todas sus facetas. El registro de contratación pública sabe quién firma contratos con el Estado, pero ignora quién montó un negocio o quién publica artículos. El registro mercantil conoce al matriculado en el RUES, pero no al investigador ni al contratista. La plataforma de ciencia y tecnología documenta al académico, pero permanece ciega ante la vida empresarial o la prestación de servicios al sector público. Cada registro nacional es una ventana parcial: nítida en su dominio, muda en los demás. De ahí nace el valor de \emph{cruzar}: solo al superponer esas ventanas parciales aparece una figura completa del graduado, y solo entonces puede afirmarse —con evidencia y no con encuestas— qué hace la USTA con quienes forma.

\subsection{El punto de partida heredado}

Este artículo no parte de cero. Continúa, con respeto y con espíritu autocrítico, el trabajo del equipo que nos precedió (Natalia González, Ximena Arias y Paula Gordillo), cuyo informe dejó una conclusión tan honesta como incómoda. Aquel equipo recolectó y catalogó con rigor las fuentes que alimentarían el observatorio —CvLAC y GrupLAC de Minciencias, el RUES de las cámaras de comercio y el SNIES del Ministerio de Educación— y describió bien cada una por separado. Sin embargo, reconoció que faltaba lo esencial: no existió un identificador que enlazara a un mismo egresado entre fuentes distintas. Sin ese puente, el cruce —que es la razón de ser misma de un observatorio de participación— quedó pendiente. El informe lo expresó sin eufemismos: el observatorio había quedado \enquote{en fase de cimientos}.

Tomamos esa frase como punto de partida y no como reproche. Unos buenos cimientos son condición necesaria para levantar un edificio; lo que faltaba era levantarlo. La pregunta que organiza este artículo es directa: ¿qué se ve cuando, por fin, las fuentes se cruzan entre sí?

\subsection{Qué cambia en esta fase}

Tres decisiones metodológicas separan esta fase de la anterior. La primera es trabajar sobre un universo concreto de graduados de la USTA, no sobre muestras ni aproximaciones. La integración de un padrón institucional ampliado triplicó la base de análisis: ahora son 174.685 graduados con títulos otorgados entre 1970 y 2026, distribuidos en seis sedes encabezadas por la Sede Principal de Bogotá y acompañadas por Bucaramanga, Tunja, Villavicencio, Medellín y la modalidad a distancia de la VUAD. Cada persona que figura en ese padrón es buscada, una por una, en los registros nacionales.

La segunda decisión es resolver el enlace por documento contra el universo completo de cada registro —no contra un extracto—, y hacerlo del lado del servidor, de modo que la coincidencia se establezca sobre el identificador único de la persona y no sobre atributos frágiles. Así se cruza a los 174.685 egresados contra la totalidad de SECOP y del RUES, evitando las pérdidas y los falsos positivos propios de los cotejos manuales o por aproximación de nombres.

La tercera atiende el caso particular de CvLAC, donde el documento de identidad no está disponible. Allí la coincidencia por nombre sería peligrosamente endeble, de manera que sustituimos esa coincidencia por una validación basada en evidencia académica —la formación declarada y la pertenencia a grupos de investigación—, exigiendo respaldo verificable antes de dar por enlazada a una persona. El resultado es un cruce que ya no descansa sobre suposiciones, sino sobre identificadores y evidencia.

\subsection{Panorama de hallazgos}

El alcance del ejercicio puede resumirse en pocas cifras macro. El universo analizado comprende 174.685 graduados de la USTA, con títulos otorgados entre 1970 y 2026, distribuidos en las seis sedes ya mencionadas. Al cruzar ese universo contra los tres registros, encontramos que el 40,3\% de los graduados —70.377 personas— deja huella verificable en al menos una de las tres dimensiones de participación consideradas. Dicho de otro modo, cerca de cuatro de cada diez egresados santotomasinos aparece, hoy y de forma rastreable, como contratista del Estado, como matriculado en el RUES formal o como investigador reconocido.

Por dimensión, las tasas globales dibujan un perfil claro. La contratación pública (SECOP) es la huella más extendida, con cerca del 24\% de los graduados registrados como proveedores del Estado: 42.889 personas. La matrícula mercantil formal (RUES) alcanza alrededor del 21,5\%, con 37.544 matriculados en el RUES. La investigación (CvLAC), por su naturaleza más selectiva y por el criterio estricto de validación que aplicamos, identifica a un grupo reducido pero cualitativamente significativo de 2.547 graduados. En el extremo más exigente del cruce —la intersección de las tres dimensiones a la vez— encontramos a 227 personas que son, simultáneamente, contratistas, empresarias e investigadoras: un núcleo pequeño de egresados de participación múltiple que solo se vuelve visible cuando las fuentes se miran juntas.

\begin{cajahallazgo}
El cruce que el observatorio tenía pendiente —su deuda declarada, su \enquote{fase de cimientos}— es hoy su primer hallazgo. Al enlazar a los 174.685 graduados de la USTA contra SECOP, RUES y CvLAC, comprobamos que más de un tercio de los egresados (40,3\%) deja una huella pública y verificable de su participación profesional, empresarial o científica. El observatorio dejó de ser una promesa de datos bien guardados para convertirse en una afirmación con evidencia.
\end{cajahallazgo}

El resto del artículo recorre este hallazgo de lo general a lo particular. Primero caracterizamos al universo de graduados por sede, programa y cohorte, para saber de quiénes hablamos. Luego examinamos cada dimensión por separado —la huella en contratación pública, el tejido matriculado en el RUES y el perfil investigador—, atendiendo a su magnitud, su composición y sus matices. A continuación abordamos el cruce integrado: cómo se combinan las dimensiones, dónde se solapan y qué perfiles de participación múltiple emergen. Después proyectamos los resultados sobre el territorio, y cerramos con una discusión que interpreta los hallazgos, reconoce los límites del enlace y traza las siguientes obras sobre los cimientos ya levantados.
